\section{Conclusiones generales}

% + CONFIRMADO:
%  - Caracteristicas del video afectan a tiempo de codificacion y energia requerida para llevar a cabo proceso de codificacion
%  - Bitrate afecta al consumo general de energia y al tiempo de codificacion
%  - No obstante, al tratar de analizar los resultados de consumo en base a los valores de SI y TI, con videos de 10 segundo de duracion, se han encontrado dificultades a la hora de poder establer una correlacion debido a la similitud de los resultados de consumode cada muestra. En cambio al llevar a cabo el analisis de consumo energetico y el tiempo de codificacion sobre videos de 30 segundos de duración se observan patrones de comportamiento en los video en relacion al SI y al TI de forma mucho mas clara.
% 
%  - x264 y x265 presentan mejoras con 3 a 6 cores y estancamiento por encima de 6. 
%  - h264 y hevc presentan mejora a partir de 3 pero saturacion por encima de 3. 
% 
%  - HEVC:
%    - Valores elevaos de TI afectan al consumo energetico y al tiempo de codifcacion, mientras que mayoria de SI mayor tiempo de codificacion
%    - 
% - H264:
%    - Se aprecia que el TI afecta al consumo energetico y el SI al tiempo de codificacion. En concreto metricas que mas consumen son HSI_LTI y HSI_HTI
%    - Falta de informacion en el bit afecta en video de larga duracion LSI_LTI
% 
% - libx264.
%    -  Merricas que mas consumen son las que tiene LSI y alto HTI junto con metricas con elevado valor de HSI y HTI con videos de 10 segundo
%    - En cuanto al tiempo de codificacion es mayor cuando el codificador le falta informacion(LSI_LTI) con video de 30 segundoscd 
% 
% - libx265:
%    - Con video de 10 segundo de duracion, videos con mayor consumo son LSI_HTI y con bajo nivel de informacion (LSI_LTI)
%    - Pero en pruebas con videos de 30 segundos se observa que TI influye en tiempo de codificacion
%    - 
% 
% + PARCIAL:
% - Codificadores de CPU presentan mejora inmediatas con 3 cores pero estancamiento con 6, pero no es visible en todas las muestras. Es necesario experimentar con un procesador que cuente con mas cores. 
% 
% 
% + No CONFIRMADO.
%  - Que grupo afecta al consumo energetico, valores muy iregulares y similares entre si. 
%  - 
% 

El conjunto de experimentos realizados ha permitido analizar el consumo energético asociado a cargas de trabajo de codificación de vídeo desplegadas en contenedores, considerando diferentes herramientas de monitorización, características del contenido, configuraciones de codificación y restricciones sobre los recursos de procesamiento disponibles. En particular, se ha estudiado la influencia del \textit{bitrate}, las métricas \textit{Spatial Information} (SI) y \textit{Temporal Information} (TI), el número de \textit{threads} utilizados por el codificador y el número de \textit{cores} disponibles para el contenedor. Para ello, se han evaluado cuatro codificadores, dos basados exclusivamente en CPU y dos mediante aceleración hardware sobre GPU.

Los resultados obtenidos en el primer experimento, orientado a la selección de una herramienta adecuada para la recopilación de métricas de consumo energético, han puesto de manifiesto diferencias tanto en la precisión de las mediciones como en la capacidad de atribuir el consumo a procesos concretos. En particular, se han identificado discrepancias en las estimaciones proporcionadas por \textit{Scaphandre} y \textit{Kepler}, especialmente en escenarios en los que varios procesos comparten los recursos del nodo. Considerando conjuntamente la precisión obtenida, la granularidad de las mediciones y las funcionalidades disponibles, se seleccionó \textbf{Container Energy Meter} como la herramienta más adecuada para las fases experimentales posteriores.

En relación con los parámetros de codificación, los resultados muestran una influencia consistente del \textit{bitrate} sobre el consumo energético y el tiempo de ejecución. Independientemente del codificador utilizado y de la configuración de acceso a los recursos hardware, las cargas configuradas con unbitrate superior presentan, en general, un mayor consumo energético y un mayor tiempo de codificación. Por tanto, la selección del bitrate constituye un factor relevante para la eficiencia de las cargas de trabajo, aunque debe realizarse conjuntamente con los requisitos de calidad del vídeo y las características del códec empleado.

Por otra parte, el análisis de las métricas SI y TI ha permitido determinar que las características espaciales y temporales del contenido influyen en el comportamiento energético y temporal de los codificadores, aunque esta relación no es suficientemente estable como para establecer una regla general basada exclusivamente en dichas métricas. Las diferencias observadas entre los experimentos realizados con vídeos de distinta duración ponen de manifiesto, además, la importancia de las condiciones de ejecución a la hora de analizar esta relación. El análisis de las cargas de mayor duración ha permitido identificar tendencias más claras, especialmente en determinados codificadores, aunque se mantiene una elevada variabilidad entre muestras con valores similares de SI y TI.

En los codificadores acelerados mediante GPU se ha observado una mayor influencia de las características espaciales y temporales del contenido sobre el comportamiento de la carga. En particular, \textit{H.264} presenta una mayor sensibilidad respecto al SI en el consumo energético y al TI en determinadas configuraciones, mientras que \textit{HEVC} muestra una influencia más destacada del SI, especialmente sobre el tiempo de codificación. Estos resultados refuerzan la hipótesis de que la complejidad del contenido constituye un factor relevante en el coste de la codificación, pero también evidencian que SI y TI no son suficientes para explicar por completo las variaciones observadas.

En los codificadores basados en CPU también se ha identificado una influencia de las características del contenido. En particular, \textit{x265} presenta una sensibilidad destacada respecto al TI, tanto en términos de consumo energético como de tiempo de ejecución, mientras que en \textit{x264} las muestras caracterizadas por valores bajos de SI y TI presentan, en determinados casos, incrementos significativos del consumo y del tiempo de codificación. En consecuencia, los resultados muestran que la respuesta energética y temporal ante una determinada configuración de recursos depende no solo del codificador empleado, sino también de las características concretas del vídeo procesado.

En cuanto a la asignación de recursos de procesamiento, se han observado diferencias significativas entre los codificadores basados en CPU y aquellos acelerados mediante GPU. En los codificadores GPU, la asignación de \textit{3 cores} y \textit{3 threads} permite obtener una mejora significativa respecto a configuraciones con un único core, identificándose a partir de este punto un comportamiento próximo a la saturación tanto para \textit{H.264} como para \textit{HEVC}. En los codificadores basados en CPU, las mejoras se mantienen hasta configuraciones de entre \textit{3} y \textit{6 cores}, identificándose \textit{6 cores} como un punto a partir del cual la incorporación de recursos adicionales deja de proporcionar beneficios proporcionales.

Estos resultados permiten establecer una conclusión especialmente relevante respecto a la relación entre tiempo de ejecución y consumo energético: una reducción del tiempo necesario para completar una carga de trabajo no implica necesariamente una reducción equivalente de su consumo energético. El incremento del número de recursos puede reducir significativamente el tiempo de ejecución, pero también aumentar el consumo instantáneo asociado a la carga. Por tanto, la configuración óptima debe considerar conjuntamente el consumo energético y el tiempo de ejecución, evitando asumir que la utilización del máximo número de recursos disponibles constituye necesariamente la alternativa más eficiente.

Finalmente, la comparación entre los cuatro codificadores pone de manifiesto una ventaja significativa de las soluciones aceleradas mediante GPU frente a las basadas exclusivamente en CPU, tanto en términos de consumo energético como de tiempo de ejecución. Esta diferencia resulta especialmente relevante en el caso de \textit{x265}, que presenta los mayores costes energéticos y temporales entre las configuraciones analizadas. Dentro de las soluciones aceleradas mediante GPU, \textit{H.264} presenta una ligera ventaja frente a \textit{HEVC} en las condiciones experimentales estudiadas, tanto en consumo energético como en tiempo de ejecución. Por tanto, para escenarios con grandes volúmenes de vídeo o con múltiples cargas de codificación concurrentes, \textit{H.264} constituye, bajo las condiciones evaluadas, una alternativa especialmente eficiente.

En conjunto, los resultados obtenidos permiten concluir que no existe una configuración universalmente óptima para minimizar el coste energético de la codificación de vídeo. La eficiencia depende de la interacción entre las características del contenido, el \textit{bitrate}, el codificador seleccionado y los recursos de procesamiento asignados. En consecuencia, las hipótesis iniciales quedan confirmadas únicamente de forma parcial: el incremento de recursos puede mejorar significativamente el tiempo de ejecución, pero no garantiza una reducción del consumo energético, mientras que las métricas SI y TI permiten caracterizar parcialmente la complejidad del contenido, pero no son suficientes por sí solas para predecir su coste energético. Esta dependencia entre contenido, configuración y recursos constituye uno de los principales resultados derivados del estudio.


\section{Limitaciones}
% - Que no se ha estudiado
% - Que s puede mejorar
% - Como se podria ampliar o extender este trabajo de investigacion

% MEJORAS
% - Ampliar el numero de muestras 
% - Ampliar las metricas subjetivas y rankear aquellas que mas afectan al consumo y con que parametros
% - No se ha escogido mas metricas debido a la falta de tiempo respecto al plazo de entrega del TFM. 
% - Estudiar opcionde unikernels y comparar el consumo energetico y tiempode ejecucion respecto a los contenedores
A pesar de los resultados obtenidos, el estudio presenta una serie de limitaciones que deben tenerse en cuenta a la hora de interpretar y generalizar las conclusiones.

Una de las principales limitaciones está relacionada con la duración de las cargas de trabajo utilizadas durante los experimentos. Los codificadores acelerados mediante GPU presentan tiempos de ejecución reducidos, situándose en torno a los \textit{8--9 segundos} de media en las pruebas realizadas. Esta duración limita la capacidad para observar la evolución del consumo energético durante cargas prolongadas y dificulta determinar si el comportamiento observado se mantiene estable cuando el sistema permanece ejecutando la misma carga durante periodos más extensos. En particular, no ha sido posible analizar con suficiente profundidad posibles efectos relacionados con la estabilización de la carga, el calentamiento de los componentes o las variaciones en su utilización a lo largo del tiempo.

Otra limitación está relacionada con el número de muestras disponibles. Aunque se ha realizado una selección de vídeos con diferentes características y se han establecido diferentes configuraciones experimentales, un número mayor de muestras permitiría aumentar la significación estadística de los resultados y determinar con mayor fiabilidad si las tendencias observadas pueden generalizarse a otros contenidos, configuraciones de hardware y condiciones de ejecución.

Asimismo, el análisis de las características del contenido se ha centrado principalmente en las métricas SI y TI. Si bien estas métricas permiten obtener una aproximación a la complejidad espacial y temporal de los vídeos, los resultados muestran que no son suficientes para explicar completamente las diferencias observadas en el consumo energético. La incorporación de otras características del contenido podría proporcionar una representación más completa de la complejidad de las cargas de codificación.

El estudio también está condicionado por las características del hardware disponible para la realización de los experimentos. En particular, las pruebas con aceleración hardware se han realizado utilizando una NVIDIA Quadro M4000. Por tanto, los resultados obtenidos no pueden extrapolarse directamente a otras arquitecturas de GPU, generaciones de hardware o configuraciones con diferentes capacidades de codificación. Del mismo modo, las conclusiones obtenidas para los codificadores utilizados no necesariamente son aplicables a otros codificadores o estándares de vídeo que no hayan sido evaluados.

Otra limitación se encuentra en el número de configuraciones de codificación analizadas. El espacio de parámetros disponible en \textbf{FFmpeg} es considerablemente amplio y el estudio se ha limitado a un subconjunto de configuraciones, seleccionadas en función de los objetivos y recursos disponibles para el proyecto. Por tanto, podrían existir otras combinaciones de parámetros capaces de proporcionar una relación diferente entre calidad, tiempo de ejecución y consumo energético.

Finalmente, las herramientas empleadas para la monitorización presentan sus propias limitaciones de precisión, granularidad y atribución del consumo energético. La estimación del consumo de un proceso dentro de un nodo que ejecuta simultáneamente otras cargas de trabajo constituye un problema complejo, por lo que los valores obtenidos deben interpretarse teniendo en cuenta las características y el margen de error asociado a cada mecanismo de medición.

Estas limitaciones no invalidan los resultados obtenidos, pero sí delimitan el alcance de las conclusiones y ponen de manifiesto la necesidad de realizar estudios adicionales con cargas de mayor duración, un conjunto más amplio de vídeos, diferentes arquitecturas de hardware y un mayor número de parámetros de codificación. La ampliación del número de muestras y variables analizadas permitiría, además, determinar con mayor precisión los factores que condicionan el consumo energético y avanzar hacia modelos capaces de predecir y optimizar automáticamente dicho consumo.

\section{Ampliaciones y Trabajo Futuro}
% - Desarrollo de ia
% - Clasificacion masiva de caracteristicas
% - Estudio metricas de consumo utilizando unikernels frente a docker container
% - Estudio presets y configuracion fina de ffmpeg
% - Ampliacion de estudio de limitacionde numero de cores con mayor cantidad de metricas mas variadas y con valores de SI TI mas elevados.